\documentclass[10pt]{article}
\usepackage[margin=1.8cm]{geometry}
\usepackage{graphicx}
\usepackage{amsmath}
\usepackage[hidelinks]{hyperref}
\setlength{\parskip}{2pt}\setlength{\parindent}{0pt}

\title{\vspace{-1.4cm}Neuron mechanisms I: how $d2$-dominant neurons compute pair parities\\
\large \texttt{wd3\_k0\_N2048}, step 500k --- post-nonlinearity units, uniform evaluation data}
\author{}\date{\vspace{-1.2cm}}

\begin{document}\maketitle\vspace{-0.9cm}

\section*{Setup and mechanism}
For each neuron we analyze $a=\phi(z)$, with $\phi=\mathrm{relu}(\tanh)$ and
$z=(Vh)_n$, evaluated on all $2^{16}$ uniform inputs. We call a neuron
\emph{$d2$-dominant} when its largest coefficient among the tracked tree modes is a
degree-2 pair parity $\chi_{x_ix_{i+1}}$. Such neurons make up most of the active
population in every block (1528/1330/1061/1036 active per block).

The exemplar (block~2, neuron~1779, pair $d2\_4=(x_8,x_9)$) makes the mechanism
particularly clear (Fig.~\ref{fig:ex}). Its pre-activation is almost exactly a
\emph{balanced difference of the two singleton inputs},
$z\approx\gamma(x_8-x_9)$ with $\gamma=0.173$ and residual std $0.004$. Therefore
$a=\phi(z)$ is positive only at the single corner
$(x_8,x_9)=({+},{-})$, where its value is
$\tanh(2\gamma)\approx0.35$.

A detector for one corner of the two-bit hypercube has the Walsh decomposition
$a=\frac{t}{4}(1+x_8-x_9-x_8x_9)$. Thus it contains \textbf{four Walsh coefficients
of equal magnitude}: a constant term, the two singleton modes, and the pair parity.
Since the constant term does not contribute to activation variance, the variance is
split into exact thirds over $\{x_8,\,x_9,\,x_8x_9\}$ (measured
$0.330/0.330/0.327$).

The target tree does not include degree-1 modes, so only the $x_8x_9$ component is
counted as tracked tree power. This is why the tree coverage is $1/3$
\emph{exactly by the structure of the corner detector}: the remaining two thirds
are the two singleton components $x_8$ and $x_9$, each with essentially the same
coefficient magnitude as the pair itself.

\begin{figure}[h!]\centering
\includegraphics[width=0.74\textwidth]{fig0_d2_exemplar.pdf}\vspace{-2pt}
\caption{The $d2$ corner unit (block 2, \#1779). Left: it fires on one corner of
$(x_8,x_9)$. Middle: measured Walsh spectrum (bars) vs.\ the parameter-free ideal
corner unit (ticks) --- exact thirds; only the pair (blue) is tree-tracked. Right:
the pre-activation is a balanced singleton difference.}
\label{fig:ex}
\end{figure}

\section*{Population: every active $d2$ neuron is a corner unit}
The exemplar is not an isolated case. Rank the $d2$-dominant neurons in each block
by their pair power $a^2$. Among all neurons before the dead/noise tail, the fraction
whose pre-activation is a balanced two-singleton combination is \textbf{1.0 in every
block}. Here ``balanced'' means that the two relevant singletons account for at least
$90\%$ of $\mathrm{Var}\,z$ and their coefficient magnitudes satisfy
$|\alpha|/|\beta|\geq0.8$. For the top 256 neurons in each block, the median fraction
of pre-activation variance coming from the two relevant singletons is
$0.993$--$0.999$ (Fig.~\ref{fig:pop}).

All four possible corner-detector sign patterns are represented. Roughly half of the
neurons are difference type,
$z\propto x_i-x_{i+1}$, which fires on a corner with pair parity $-1$ and therefore
has a negative pair coefficient. The other half are sum type,
$z\propto x_i+x_{i+1}$, which fires on a corner with pair parity $+1$ and therefore
has a positive pair coefficient.

This balanced population is exactly what allows the network to reconstruct the full
$\pm1$ pair parity downstream. Individual neurons carry the desired pair mode together
with singleton by-products, but an appropriate weighted sum across many corner units
can reinforce the pair component while cancelling the singleton components.

Consistent with the $d4$/$d8$ reports, these pair features are \emph{end products}
for the supervised $d2$ targets rather than intermediate building blocks for the
higher-degree parities. The deeper parity modes are generated by separate raw-bit
majority-gate mechanisms rather than by composing these $d2$ corner units.

\begin{figure}[h!]\centering
\includegraphics[width=0.9\textwidth]{d2_corner_fraction_vs_rank.pdf}\vspace{-2pt}
\caption{Fraction of $d2$-dominant neurons with corner input vs.\ rank by pair power
(blue, pinned at 1.0 until the noise tail; gray: $a^2$ of the pair); green: the strict
difference type $z\propto x_i-x_{i+1}$ ($\approx\!1/2$).}
\label{fig:pop}
\end{figure}

\vspace{2pt}\footnotesize
Artifacts: \texttt{gdrive:parity/analysis\_output/neuron\_mechanisms/wd3\_k0\_N2048/step\_00500000/}
(\texttt{per\_neuron.csv}, figures, this report; companion $d4$/$d8$/$d16$ reports in
the same folder).
\end{document}
